\section{Formal Sketch}
\label{sec:formal}

The verifier of \S\ref{sec:predicate} admits a machine-checkable transcription of its accept relation in Lean 4. The transcription is not a deep theorem and is not load-bearing for the cryptographic security claim; it is an audit artifact that pins the informal predicate to a mechanical object so that the schema and the verifier cannot silently drift apart. Strip the runtime to three structures: a trust store $\mathcal{S}$ with separate issuer-key and kernel-key allowlists, a bilateral envelope $E$ carrying a receipt id, a scope predicate, and two signature slices (one issuer-side, one kernel-side), and a denotational interpreter $\llbracket\cdot\rrbracket$ for the scope predicate against the receipt id. The accept relation is the conjunction
\begin{align*}
  \mathit{accept}(\mathcal{S},E) \,=\;
    & [E.\mathit{issuerSig}.\mathit{kid} \in \mathcal{S}.\mathit{issuerKeys}] \\
    {}\wedge{} & [E.\mathit{kernelSig}.\mathit{kid} \in \mathcal{S}.\mathit{kernelKeys}] \\
    {}\wedge{} & \llbracket E.\mathit{scope}\rrbracket(E.\mathit{receiptId}).
\end{align*}
Signature byte validity is abstracted into trust-store membership: the trust store names keys whose signatures the deployed verifier has separately checked, and the Lean transcription inspects only the membership and scope conjuncts. The cryptographic primitives sit underneath this abstraction and are out of the Lean module's scope.

\paragraph{Three abstract gates vs.\ six runtime gates.} The six runtime gates of \S\ref{sec:predicate} (canonical-payload, predicate-type, signer-distinctness, trust-store-membership, lease-freshness, subject-digest) collapse into the three structural conjuncts above. Four of the runtime gates -- canonical-payload, predicate-type, signer-distinctness, and lease-freshness -- are assumed-sound at the byte-canonicalization and verifier-state layer that the trust-store conjuncts sit above: the trust store names keys whose signatures, predicate-type, and signer-distinctness checks the deployed verifier has separately discharged, and the lease-freshness check is a verifier-state precondition that yields to the trust-store conjuncts as a wrapper. The two load-bearing structural gates that remain visible to Lean are issuer-key membership and kernel-key membership, plus the scope-predicate denotation that the subject-digest gate witnesses at runtime. The transcription is therefore a structural projection, not a runtime equivalence.

\begin{theorem}[\thm{freestanding_accept_set_theorem}]
For every trust store $\mathcal{S}$ and bilateral envelope $E$,
\begin{align*}
  \mathit{accept}(\mathcal{S},E) \;\Longleftrightarrow\;
    & E.\mathit{issuerSig}.\mathit{kid} \in \mathcal{S}.\mathit{issuerKeys} \\
    {}\wedge{} & E.\mathit{kernelSig}.\mathit{kid} \in \mathcal{S}.\mathit{kernelKeys} \\
    {}\wedge{} & \llbracket E.\mathit{scope}\rrbracket(E.\mathit{receiptId}) = \mathit{true}.
\end{align*}
\end{theorem}

The bi-conditional is definitional: $\mathit{accept}$ is defined as exactly that Boolean conjunction lifted through \codepath{decide}, and the proof is one Lean tactic line that unfolds the definition and rewrites the resulting Boolean form. The statement therefore witnesses that the verifier's accept set decomposes into three abstract conjuncts -- two trust-store-membership tests and one scope denotation -- under the abstraction discussed above, with the canonical-payload, predicate-type, signer-distinctness, and lease-freshness runtime gates of \S\ref{sec:predicate} treated as preconditions of the trust-store conjuncts. Its work is to foreclose on accidental drift between the informal schema and the mechanical verifier as the predicate evolves; it does not constitute an independent proof that the cryptographic gates compose into a sound admission relation. The Lean source is self-contained at the module level: it imports the predicate-language type and its denotational interpreter, but takes no dependency on the wider polity model, anchor-quorum machinery, or treaty-intersection construction. A reviewer can audit the file end-to-end at \codepath{formal/lean4/Chio/Chio/Treaty/BilateralAccept.lean}.

Three corollaries follow definitionally and serve as scope-clarification artifacts rather than as independent theorems. The first, \thm{accept_monotone_in_issuer_store}, names accept-set monotonicity in the trust store: if every issuer key in $\mathcal{S}$ is also in $\mathcal{S}'$ and the kernel-key allowlists agree, then $\mathit{accept}(\mathcal{S},E) = \mathit{true}$ implies $\mathit{accept}(\mathcal{S}',E) = \mathit{true}$. The statement is the List-membership-monotonicity lemma over the issuer-key conjunct; its purpose is to record explicitly that adding a co-signing counterparty's key to the issuer allowlist never causes a previously-admitted envelope to be rejected. The second, \thm{accept_conj_scope_decompose}, names accept-set conjunction over scope predicates: an envelope whose scope is $p \wedge q$ and which is accepted under $\mathcal{S}$ yields the two trust-store conjuncts plus $\llbracket p\rrbracket(E.\mathit{receiptId}) = \mathit{true}$ and $\llbracket q\rrbracket(E.\mathit{receiptId}) = \mathit{true}$. This is the denotational identity $\llbracket p \wedge q\rrbracket = \llbracket p\rrbracket \wedge \llbracket q\rrbracket$ stated without reusing the original signatures for a byte-distinct rebound envelope. The third, \thm{accept_requires_issuer_key}, is the contrapositive of the first conjunct: if the envelope's issuer key id is not in $\mathcal{S}.\mathit{issuerKeys}$, then $\mathit{accept}(\mathcal{S},E) = \mathit{false}$. All three are proved by unfolding $\mathit{accept}$ and rewriting on the main bi-conditional.

\paragraph{Kernel-only axioms.} Running Lean's \codepath{\#print axioms} against \thm{freestanding_accept_set_theorem} and each of the three corollaries reports only the kernel-level axioms \emph{propext}, \emph{Classical.choice}, and \emph{Quot.sound}: the same trusted base every Mathlib development presupposes. No \emph{sorry} placeholders appear, and the module introduces no custom axioms. This is a hygiene claim, not a security claim: every Lean development that avoids \emph{sorry} and \emph{axiom} declarations enjoys the same property, and the audit value is that any future strengthening of the module can be detected by an axiom-footprint diff against this baseline.

\paragraph{What the Lean module does not prove.} The transcription is bounded by four explicit gaps. First, Ed25519 byte-level signature soundness: the verifier's signature check is an axiomatic primitive whose correctness lives outside the module, and the trust-store membership conjuncts presuppose that the deployed verifier has separately discharged it. Second, SHA-256 collision resistance and JCS injectivity: the subject digest is treated as a hash over canonical bytes, and any cross-field collision or canonicalization disagreement at the cryptographic layer is a structural assumption rather than a Lean theorem. Third, Rust-Lean refinement: the production verifier in \codepath{verify_chio_bilateral_dsse_envelope} is not extracted from this module nor proven equivalent to it, and the two artifacts share no formal connection beyond the informal schema in \S\ref{sec:predicate}. Fourth, polity-level admission: the module speaks only of a single bilateral envelope and does not address treaty intersection, single-amendment backward refinement, or the trajectory-invariant rule that closes constitutional ratchets; \thm{amendment_admissible_iff_backward_refinement}, \thm{treaty_admission_iff_predicate_intersection}, and \thm{essential_preserved_chain} are discharged in the broader polity-level formalization (anonymized for review). The bilateral-DSSE primitive sits below those obligations: it is the byte-level admission gate the polity model presupposes, and the Lean module's job is to keep the gate's mechanical statement aligned with the schema rather than to derive the polity model's properties.
